\section{A general bound for more than three sectors}\label{four:section}
The same comparison and witness arguments extend to any fixed number of
sectors. The resulting bound is sharp when there are three sectors, by the
main theorem, but no sharpness assertion is made for more than three.
This extension concerns the same pointwise protection of original core
maxima; it imposes no strategic reporting or incentive requirement.

Let $K\ge3$ and let $N_1,\ldots,N_K$ be disjoint sectors, each of size
$n\ge1$. For each unordered pair of distinct sectors fix a nonnegative
real pair-value matrix. A productive full team contains one player from
every sector, and its worth is the sum of all its pair values. Coalition
worth is the largest total worth of player-disjoint full teams it contains.
All smaller teams have worth zero. Assume $T=v(N)>0$ and that the entire
original core $\core$ is nonempty. Define $H_i=\max_{x\in\core}x_i$ and
\[
 s_K(M)=\min\left\{\sum_i y_i-T:
       y_i\ge H_i/K\ (i\in N),\quad y(e)\ge v(e)
       \text{ for every full team }e\right\}.
\]
As in the three-sector model, these constraints imply nonnegativity and
coverage of every original coalition.

\begin{theorem}[General sector-count bound]\label{four:main}
Every finite market in the preceding class satisfies
\begin{equation}\label{four:universal}
             s_K(M)<\frac{2(K-2)}{K(K-1)}\,T.
\end{equation}
For an integral partial matching $P$ leaving $k=n-|P|\ge1$ unmatched
players in each sector, put
\[
 J_K(M,P)=v(P)+\frac1K\sum_{i\notin\cup P}H_i-T.
\]
Then
\begin{equation}\label{four:integral}
 J_K(M,P)\le
       \frac{2(K-2)(k-1)}{K\bigl(k(K-1)-1\bigr)}\,T.
\end{equation}
If $k=0$, then $J_K(M,P)\le0$.
\end{theorem}

We give the proof because the fractional-packing step and the use of
separate witnesses are essential. Choose an optimal integral matching,
complete it to a perfect matching, and relabel its full teams by
$j=0,\ldots,n-1$. Completion preserves optimality: the sector sizes are
equal, and the added full teams have nonnegative worth. Write $d_j$ for
the worth of diagonal team $j$, and $u_{a,j}$ for the sum of its pair
values incident with its sector-$a$ player. Thus
\begin{equation}\label{four:budgets}
       \sum_j d_j=T,\qquad \sum_{a,j}u_{a,j}=2T.
\end{equation}
Every diagonal team is tight at every $x\in\core$. In particular,
\[
   \sum_{a=1}^Kx_{a,j}=d_j,\qquad
   0\le H_{a,j}\le d_j,\qquad
   \delta_{a,j}=d_j-H_{a,j}\ge0.
\]
Replacing player $(a,j)$ in its diagonal team by $(a,h)$ leaves all
pairs not incident with sector $a$ intact. Nonnegativity of the new
pairs therefore gives
\begin{equation}\label{four:range}
             x_{a,j}-x_{a,h}\le u_{a,j}
             \qquad(x\in\core).
\end{equation}
There is no division by $d_j$, so zero diagonal values are allowed.

\begin{lemma}[Aggregate comparisons for $K$ sectors]\label{four:aggregate}
Let $\lambda$ be any fractional partial matching, with
\[
 p_i=\sum_{e\ni i}\lambda_e\le1,\quad q_i=1-p_i,\quad
 W=\sum_e\lambda_ev(e),\quad S=\sum_iq_iH_i,\quad V_K=KW+S.
\]
Set
\[
 k=n-\sum_e\lambda_e,\qquad
 \beta=\sum_iq_iu_i,\qquad C_2=\sum_iq_i^2u_i,
\]
where $u_i=u_{a,j}$ for $i=(a,j)$, and put
\[
 E=\sum_{a,j}\left(\sum_{b\ne a}q_{b,j}\right)\delta_{a,j}\ge0.
\]
Then
\begin{align}
 V_K&\le (3K-4)T-(K-2)\beta-\frac{E}{K-1},
        \label{four:first}\\
 V_K&\le KT+\beta-\frac{C_2}{k}\qquad(k>0).
        \label{four:second}
\end{align}
\end{lemma}
\begin{proof}
For a full team $e=((1,j_1),\ldots,(K,j_K))$ and distinct sectors
$b,c$, subtract diagonal tightness at $j_b$ from coverage of $e$.
For every $x\in\core$, apply~\eqref{four:range} to the other $K-2$
sectors to obtain
\[
 x_{c,j_c}-x_{c,j_b}
 \ge v(e)-d_{j_b}-\sum_{a\notin\{b,c\}}u_{a,j_a}.
\]
Choose an original core point attaining the subtracted maximum
$H_{c,j_b}$, while bounding its other coordinate by $H_{c,j_c}$.
This proves the same inequality for the difference of those two maxima.
Each ordered pair $(b,c)$ may use a different attaining core point.
Sum over all $K(K-1)$ ordered pairs, writing
$H_j^\Sigma=\sum_aH_{a,j}$:
\begin{equation}\label{four:team-comparison}
\begin{split}
 K(K-1)v(e)\le{}&(K-1)\sum_a d_{j_a}
       +(K-1)(K-2)\sum_a u_{a,j_a}\\
       &+K\sum_aH_{a,j_a}-\sum_aH_{j_a}^\Sigma.
\end{split}
\end{equation}
Here every incident budget occurs $(K-1)(K-2)$ times, and the ideal
terms are $K\sum_aH_{a,j_a}-\sum_aH_{j_a}^{\Sigma}$.
The count preserves multiplicities when some team labels coincide.

Multiply~\eqref{four:team-comparison} by $\lambda_e$ and sum over
teams. Add $(K-1)S$. If $n_j=\sum_bp_{b,j}$, the resulting coefficient
of $H_{a,j}$ is
\[
       Kp_{a,j}+(K-1)q_{a,j}-n_j
                   =\sum_{b\ne a}q_{b,j}.
\]
Replacing $H_{a,j}$ by $d_j-\delta_{a,j}$ gives an ideal contribution
$(K-1)\sum_{a,j}q_{a,j}d_j-E$. The other diagonal contribution is
$(K-1)\sum_{a,j}p_{a,j}d_j$. These sum to $K(K-1)T-E$.
The budget contribution is
$(K-1)(K-2)(2T-\beta)$ by~\eqref{four:budgets}.
The left-hand side is $(K-1)V_K$, proving~\eqref{four:first}.

Every sector has residual mass $k$: each full team uses exactly one
player per sector. For every original core point $x$, fractional team
coverage gives $\sum_i p_ix_i\ge W$, and hence
\begin{equation}\label{four:residual}
                   \sum_iq_ix_i\le T-W.
\end{equation}
In particular $W\le T$. Fix a player $i$ and choose a core point
attaining $H_i$. Its other same-sector coordinates are at least
$H_i-u_i$ by~\eqref{four:range}, while all other-sector coordinates
are nonnegative. Therefore
\[
        T-W\ge q_iH_i+(k-q_i)(H_i-u_i),
        \qquad
        kH_i\le T-W+(k-q_i)u_i.
\]
The factor $k-q_i$ is the sum of the other residual weights in $i$'s
sector and is nonnegative. Multiplying by $q_i$, summing the resulting
scalar inequalities, and using $\sum_iq_i=Kk$ yields
\[
                    kS\le Kk(T-W)+k\beta-C_2.
\]
For $k>0$, division by $k$ proves~\eqref{four:second}. Separate
maximizing witnesses suffice because only the fixed numbers $H_i$
appear in these scalar inequalities.
\end{proof}

\begin{proof}[Proof of Theorem~\ref{four:main}]
The shifted covering-program dual is
\begin{equation}\label{four:dual}
 s_K(M)=\max_{\lambda\in\pack_K}\left(W+\frac SK-T\right)
       =\max_{\lambda\in\pack_K}\left(\frac{V_K}{K}-T\right),
\end{equation}
where $\pack_K$ is the fractional partial-matching polytope just defined.
Indeed, substitute $y_i=H_i/K+z_i$, $z_i\ge0$, in the covering program;
its dual has the load constraints $p_i\le1$. Adding the constant
$\sum_iH_i/K-T$ gives~\eqref{four:dual}. The primal is feasible and
bounded below, while $\pack_K$ is compact, so both optima are attained.

For $k>0$, take weights $1/(K-1)$ and $(K-2)/(K-1)$ on
\eqref{four:first} and~\eqref{four:second}, respectively. The
$\beta$ terms cancel, giving
\begin{equation}\label{four:fractional-combination}
 V_K\le\left(K+\frac{2(K-2)}{K-1}\right)T
         -\frac{E}{(K-1)^2}
         -\frac{(K-2)C_2}{(K-1)k}.
\end{equation}
If $C_2>0$, this is strictly smaller than the displayed leading term.
If $C_2=0$, nonnegativity of each $q_i^2u_i$ implies $\beta=0$, so
\eqref{four:second} gives $V_K\le KT$, again strictly smaller since
$K\ge3$ and $T>0$. When $k=0$, all residual weights vanish and
\eqref{four:residual} gives $W\le T$, so $V_K\le KT$ as well.
Apply this strict bound at a maximizing point of the compact dual
polytope to prove~\eqref{four:universal}. A pointwise strict inequality
without attained maximization would not by itself justify this step.

For an integral partial matching with $k\ge1$, residual weights belong
to $\{0,1\}$ and $C_2=\beta$. Drop $E\ge0$ from
\eqref{four:first} and take weights
\[
       \frac{k-1}{k(K-1)-1},\qquad
       \frac{k(K-2)}{k(K-1)-1}
\]
on~\eqref{four:first} and~\eqref{four:second}. They are nonnegative
and sum to one. The $\beta$ terms again cancel, giving
\[
       V_K\le\left(K+
           \frac{2(K-2)(k-1)}{k(K-1)-1}\right)T.
\]
Since $J_K=V_K/K-T$, this proves~\eqref{four:integral}, including
$k=1$. For $k=0$, the matching is perfect and $J_K=v(P)-T\le0$.
\end{proof}

\begin{remark}[Four sectors and the scope of sharpness]\label{four:scope}
For $K=4$, the theorem gives $s_4(M)<T/3$ and
$J_4(M,P)\le(k-1)T/(3k-1)$. It therefore gives a supremum upper bound
$\sup_M s_4(M)/T\le1/3$; it does not prove that this supremum equals
$1/3$, nor a smaller fixed gap below $1/3$ uniform over market sizes.
The analogous distinction applies for every $K>3$. The integral bound
is an upper bound on one integral dual obstruction, not on the full
fractional optimum. When $K=3$, both formulas recover the sharp bounds
already proved in the main text. Sparse allowed sector graphs are
included by setting the unused pair matrices to zero.
\end{remark}

\begin{remark}[What the exchange budgets provide]\label{four:exchange}
More generally, the same argument works for full-team games with a tight
optimal block partition and nonnegative numbers $u_i$ satisfying the
directed same-sector range bound~\eqref{four:range}. If
$B=\sum_i u_i$, the first bound becomes
\[
       V_K\le KT+(K-2)(B-\beta)-E/(K-1),
\]
and the same weighted combination gives
$s_K(M)\le (K-2)B/[K(K-1)]$.
Pair nonnegativity supplies these range bounds through an actual
replacement-team calculation, and pair additivity supplies $B=2T$.
These hypotheses cannot be replaced merely by assumed versions of the
two aggregate inequalities. This is a direct extension of the proof,
not a separately established sharpness or historical-priority claim.
\end{remark}
